
\documentclass[12pt,a4paper]{article}

\usepackage[utf8]{inputenc}
\usepackage[T1]{fontenc}
\usepackage{graphicx}
\usepackage{array}
\usepackage{geometry}
\usepackage{placeins} % Keep diagrams and tables near their explanations

\geometry{
    top=2cm,
    bottom=2cm,
    left=2.5cm,
    right=2.5cm
}

\begin{document}

\begin{titlepage}
\centering

% ========== FPT LOGO ==========
\includegraphics[width=5cm]{../../reports/overleaf/src/logofpt.png}

\vspace{1.1cm}

{\Large\bfseries FPT UNIVERSITY\par}

\vspace{0.8cm}

{\large INTERNET OF THINGS -- IOT102\par}

\vspace{1.4cm}

{\LARGE\bfseries FINAL PROJECT REPORT\par}

\vspace{0.6cm}

{\Large\bfseries
IoT-Based Fire Detection\\
and Warning System\par}

\vspace{1.1cm}

{\large Class: SE2027 \qquad Group: G2\par}

\vspace{0.8cm}

% ========== GROUP MEMBERS ==========
\renewcommand{\arraystretch}{1.35}

\begin{tabular}{p{8cm} p{3.5cm}}
\hline
\textbf{Group Member} & \textbf{Student ID}\\
\hline
Hoang Nhat Tam & SE202018\\
Phan Tan Hung & SE201027\\
Le Vo Tuan Kiet & SE203376\\
\hline
\end{tabular}

\vspace{1.0cm}

{\large\textbf{Instructor:} Nguyen Duc Loi\par}

\vfill

{\large Ho Chi Minh City\par}
{\large November 2026\par}

\end{titlepage}

% ========== REPORT CONTENT ==========

% ==========================================
% FINAL PROJECT REPORT CONTENT
% ==========================================

\begin{abstract}
Fire incidents can cause significant damage to human life
and property when early warning mechanisms are unavailable
or ineffective. This project presents an ESP32-based fire
detection and warning prototype integrating an LM35
temperature sensor, an MQ-2 smoke/gas sensor, and a flame
sensor. The system employs a rule-based detection algorithm
to classify environmental conditions into three states:
SAFE, WARNING, and DANGER. Local feedback is provided
through an LCD1602 display, three status LEDs, and an
active buzzer. An SG90 servo motor simulates automatic
door opening during dangerous conditions, while a push
button supports manual door control when the system is
not in the DANGER state. The prototype was evaluated
using eight predefined functional test scenarios covering
sensor detection, state transitions, and door-control
behavior. All eight scenarios were reported as passing
during group testing. The current implementation operates
locally without cloud connectivity. Future development
may include remote monitoring, improved sensor
calibration, and false-alarm reduction techniques.
\end{abstract}

\noindent\textbf{Keywords:}
Internet of Things, ESP32, Fire Detection,
MQ-2, LM35, Flame Sensor, Early Warning System.

% ==========================================
\section{Introduction}
% ==========================================

Fire detection and early warning play an important role
in protecting residential environments, workplaces,
and other indoor facilities. Undetected flames, abnormal
temperature increases, and smoke or gas accumulation
may lead to hazardous conditions. Early recognition of
these indicators can support faster responses and
reduce potential risks.

Traditional fire warning systems commonly depend on
individual sensing mechanisms or local alarms.
However, environmental disturbances such as cooking
smoke may trigger unwanted warnings, while a single
sensor may not adequately represent different fire
conditions \cite{yu2023multidetector,fathoni2025multisensor}.
Recent studies have investigated multisensor detection
and more advanced classification methods to improve
fire monitoring.

This project develops a low-cost fire detection and
warning prototype using an ESP32 microcontroller.
Three sensing components are integrated: an LM35
temperature sensor, an MQ-2 smoke/gas sensor, and a
flame sensor. The ESP32 processes their readings
locally and determines one of three operating states:
SAFE, WARNING, or DANGER.

In addition to visual and audible warnings, the
system uses an SG90 servo motor to demonstrate an
automatic door-opening mechanism. A physical push
button allows manual control when no dangerous
condition is active.

The main objectives of this project are to:

\begin{enumerate}
    \item Integrate temperature, smoke/gas, and
    flame detection using an ESP32.
    
    \item Implement a three-level warning algorithm
    based on sensor signals.
    
    \item Provide local warnings through LEDs,
    an LCD, and a buzzer.
    
    \item Demonstrate automatic and manual
    door-control behavior.
    
    \item Verify system functionality through
    predefined test scenarios.
\end{enumerate}

The present implementation focuses on local
monitoring and actuator control. Wireless monitoring
and configurable warning policies are considered
future improvements rather than completed features.

% ==========================================
\subsection{Related Work}
% ==========================================

Several studies have explored multisensor approaches
to fire detection.

Yu et al. \cite{yu2023multidetector} investigated
real-time fire alarm technology using signal similarity
between multiple detectors. Their work combined
temporal signal alignment and correlation analysis
to support fire-event classification.

Fathoni et al. \cite{fathoni2025multisensor} proposed
a multisensor fire and hazardous-gas detection system
using an artificial neural network based on a Radial
Basis Function Network. Their system processed flame,
gas, temperature, and humidity information.

Purba et al. \cite{purba2026smartsafety} developed
an ESP32-based multi-hazard monitoring system
using a hierarchical decision-tree approach.
The system classified environmental conditions
into SAFE, ALERT, and DANGER states.

Ramadani \cite{ramadani2026blynk} presented an
ESP32-based early fire detection system integrating
flame and gas sensors with Blynk monitoring and
Telegram notifications. This work demonstrates
an approach to extending local detection with
remote monitoring capabilities.

Other studies have examined gas sensor calibration,
temporal signal analysis, and fuzzy-logic-based
hazard assessment
\cite{solorzano2022gas,nakip2021rtpnn,
rehman2021smartfire}.

Compared with these approaches, the proposed
G2 prototype employs a simpler rule-based algorithm.
This design reduces implementation complexity and
does not require a machine-learning training dataset.
However, its fixed thresholds and digital detection
inputs may be susceptible to environmental
interference and false alarms.

% ==========================================
\section{Methods and Materials}
% ==========================================

\subsection{System Architecture}

The proposed prototype consists of three main
functional layers: sensor acquisition, local
processing, and warning/actuation.

The sensing layer includes the LM35 temperature
sensor, MQ-2 smoke/gas sensor, and flame sensor.
These devices provide environmental information
to the ESP32 microcontroller.

The processing layer evaluates sensor inputs using
a prioritized rule-based algorithm. The result
determines the current warning state.

The output layer includes three LEDs, an active
buzzer, a 16x2 I2C LCD, and an SG90 servo motor.
The manual push button provides an additional
input for door control.

Figure~\ref{fig:block} illustrates the system
architecture. Remote monitoring components
present in the proposed architecture are not
part of the implemented firmware.

% ===== FIGURE 1: BLOCK DIAGRAM =====

\begin{figure}[!htbp]
    \centering
    \includegraphics[
        width=\textwidth,
        height=0.58\textheight,
        keepaspectratio
    ]{src/BlockDiagram_IOT102_G2-BlockDiagram.drawio.png}

    \caption{System architecture of the ESP32-based
    fire detection prototype. Cloud monitoring
    represents a proposed future extension.}
    \label{fig:block}
\end{figure}
\FloatBarrier



\subsection{Hardware Components}

The main controller is an ESP32 DevKit board
with a 30-pin configuration and USB Type-C
connection. Its GPIO interfaces support digital
sensor input, analog measurement, I2C communication,
and actuator control.

Table~\ref{tab:hardware} summarizes the
main components.

\begin{table}[!htbp]
\centering
\caption{Main Hardware Components}
\label{tab:hardware}
\renewcommand{\arraystretch}{1.2}
\begin{tabular}{|p{4cm}|p{9cm}|}
\hline
\textbf{Component} & \textbf{Function} \\
\hline
ESP32 DevKit &
Central processing and control unit. \\
\hline
LM35 &
Measures temperature through an analog voltage. \\
\hline
MQ-2 &
Provides a digital smoke/gas detection signal. \\
\hline
Flame Sensor &
Provides a digital flame detection signal. \\
\hline
LCD1602 I2C &
Displays sensor information and system status. \\
\hline
Three LEDs &
Indicate SAFE, WARNING, and DANGER. \\
\hline
Active Buzzer &
Generates audible warning signals. \\
\hline
SG90 Servo &
Simulates automatic door opening. \\
\hline
Push Button &
Allows manual door control. \\
\hline
\end{tabular}
\end{table}
\FloatBarrier

\subsection{Circuit Design and GPIO Mapping}

The prototype uses dedicated ESP32 GPIO pins
for sensor inputs and output devices.

The final firmware pin assignments are summarized
in Table~\ref{tab:gpio}.

\begin{table}[!htbp]
\centering
\caption{ESP32 GPIO Assignments}
\label{tab:gpio}
\renewcommand{\arraystretch}{1.15}
\begin{tabular}{|p{7cm}|p{3cm}|}
\hline
\textbf{Device / Signal} &
\textbf{ESP32 GPIO} \\
\hline
LM35 Analog Output & GPIO34 \\
\hline
MQ-2 Digital Output & GPIO18 \\
\hline
Flame Sensor Digital Output & GPIO19 \\
\hline
LCD I2C SDA & GPIO21 \\
\hline
LCD I2C SCL & GPIO22 \\
\hline
Active Buzzer Control & GPIO15 \\
\hline
Manual Button & GPIO33 \\
\hline
SG90 Servo Signal & GPIO32 \\
\hline
Green LED & GPIO25 \\
\hline
Yellow LED & GPIO26 \\
\hline
Red LED & GPIO27 \\
\hline
\end{tabular}
\end{table}
\FloatBarrier

The ESP32 development board receives a 5 V
input supply. In the assembled prototype,
the LM35 and active buzzer are powered from
the board's 3.3 V output. The buzzer was
operated at a reduced supply voltage to
lower the sound level during classroom
demonstrations.

The nominal LM35 supply requirements and
the operating range of the buzzer require
separate verification against their respective
component specifications. In addition,
voltage levels from 5 V sensor modules
and I2C pull-up connections must be checked
for compatibility with ESP32 GPIO inputs.

Figure~\ref{fig:circuit} presents the circuit
schematic. It should be verified against
the physical wiring before final submission.

% ===== FIGURE 2: CIRCUIT SCHEMATIC =====

\begin{figure}[!htbp]
    \centering
    \includegraphics[
        width=\textwidth,
        height=0.55\textheight,
        keepaspectratio
    ]{src/BlockDiagram_IOT102_G2-Circuit Schematic.drawio.png}

    \caption{Circuit schematic and GPIO connections
    of the ESP32 fire detection prototype.}
    \label{fig:circuit}
\end{figure}
\FloatBarrier


\subsection{Sensor Data Processing}

The LM35 generates an analog output voltage
proportional to temperature, with a nominal
scale factor of 10 mV per degree Celsius.

The temperature is calculated as:

\begin{equation}
T = \frac{V_{\mathrm{mV}}}{10}
\end{equation}

where $T$ is the estimated temperature in
degrees Celsius and $V_{\mathrm{mV}}$ is the
measured sensor output voltage in millivolts.

The ESP32 firmware uses
\texttt{analogReadMilliVolts()} to obtain voltage
readings. Ten samples are collected, with a
5 ms delay after each reading, and their
average is used to calculate temperature.

This averaging approach helps reduce random
measurement fluctuations. However, it does not
replace independent sensor calibration.
ADC calibration and noise are discussed in
the official Espressif documentation
\cite{espressif_adc_calibration}.

The MQ-2 and flame sensor modules are read
through digital output pins. Both signals
are treated as active LOW in the firmware.
Therefore, a LOW input indicates an activated
detection condition.

\subsection{Fire Detection Algorithm}

The fire detection algorithm categorizes
environmental conditions into three states:
SAFE, WARNING, and DANGER.

Two temperature thresholds are configured:

\begin{equation}
T_{\mathrm{warning}} = 40^\circ C
\end{equation}

\begin{equation}
T_{\mathrm{danger}} = 43^\circ C
\end{equation}

The DANGER condition is defined as:

\begin{equation}
D = F \lor S \lor (T \geq 43)
\end{equation}

where $F$ represents flame detection,
$S$ represents MQ-2 detection, and $T$
represents the estimated temperature
in degrees Celsius.

If $D$ is true, the system enters DANGER.

Otherwise, a temperature of at least
$40^\circ C$ activates WARNING.

If neither condition is satisfied,
the system remains in SAFE.

Table~\ref{tab:states} summarizes the
corresponding outputs.

\begin{table}[!htbp]
\centering
\caption{System Warning States}
\label{tab:states}
\renewcommand{\arraystretch}{1.2}
\begin{tabular}{|p{2.2cm}|p{5cm}|p{5cm}|}
\hline
\textbf{State} &
\textbf{Condition} &
\textbf{Output Behavior} \\
\hline
SAFE &
No danger; temperature below 40$^\circ$C &
Green LED ON; buzzer OFF \\
\hline
WARNING &
No danger; temperature at least 40$^\circ$C &
Yellow LED ON; intermittent buzzer \\
\hline
DANGER &
Flame OR MQ-2 detection OR temperature
at least 43$^\circ$C &
Red LED ON; continuous buzzer;
door forced OPEN \\
\hline
\end{tabular}
\end{table}
\FloatBarrier

The intermittent warning signal follows
a nominal 500 ms ON and 500 ms OFF pattern.
During DANGER, the buzzer remains continuously
activated.

The SG90 servo is commanded to 90 degrees
when DANGER is detected. Manual closing
is disabled while the danger condition
remains active.

When DANGER clears, the current door state
is retained until a valid manual button
action occurs.

The button is configured with
\texttt{INPUT\_PULLUP} and a 50 ms
debounce interval.

Figure~\ref{fig:flowchart} illustrates
the main decision flow.

% ===== FIGURE 3: SYSTEM FLOWCHART =====

\begin{figure}[!htbp]
    \centering
    \includegraphics[
        width=\textwidth,
        height=0.68\textheight,
        keepaspectratio
    ]{src/BlockDiagram_IOT102_G2-FlowChart.drawio.png}

    \caption{Fire detection and actuator control
    flowchart based on the final firmware.}
    \label{fig:flowchart}
\end{figure}
\FloatBarrier



\subsection{Software Implementation}

The firmware is developed in the Arduino
programming environment using the ESP32
hardware platform.

During initialization, the program
configures GPIO directions, the ADC,
I2C communication, serial communication,
and servo control.

The LCD1602 is accessed through I2C
using GPIO21 as SDA and GPIO22 as SCL.
The firmware checks the common PCF8574
addresses 0x27 and 0x3F.

The main execution loop repeatedly performs
sensor acquisition, state classification,
warning output updates, button processing,
and LCD/Serial status reporting.

Serial communication operates at
115200 baud and provides information
for observing sensor and system states.

The implemented firmware does not
include WiFi communication, MQTT,
Node-RED, Blynk, or cloud logging.

% ==========================================
\section{Results and Discussion}
% ==========================================

\subsection{Experimental Procedure}

The prototype was evaluated through
eight predefined functional test scenarios.

The objective of these tests was to verify
whether the system produced the expected
warning state and actuator behavior
under selected input conditions.

The group reported that all eight
scenarios behaved as expected.

The evaluation focused on functional
behavior rather than statistical
fire-detection performance.

\subsection{Functional Testing Results}

Table~\ref{tab:tests} summarizes the
test scenarios and their reported results.

\begin{table}[!htbp]
\centering
\caption{Functional Test Results}
\label{tab:tests}
\renewcommand{\arraystretch}{1.15}
\small
\begin{tabular}{|p{1cm}|p{5.2cm}|p{5cm}|p{1.2cm}|}
\hline
\textbf{ID} &
\textbf{Test Scenario} &
\textbf{Expected Behavior} &
\textbf{Result} \\
\hline
TC01 &
Normal temperature, no flame or gas/smoke detection &
SAFE; green LED ON &
PASS \\
\hline
TC02 &
Temperature in warning range; no other danger signal &
WARNING; yellow LED and intermittent buzzer &
PASS \\
\hline
TC03 &
Temperature at least 43$^\circ$C &
DANGER; continuous alarm &
PASS \\
\hline
TC04 &
Flame sensor activated &
DANGER; door opens &
PASS \\
\hline
TC05 &
MQ-2 digital detection activated &
DANGER; door opens &
PASS \\
\hline
TC06 &
Manual button pressed during DANGER &
Door remains open &
PASS \\
\hline
TC07 &
Danger condition removed &
Door retains its open state &
PASS \\
\hline
TC08 &
Manual button pressed during SAFE &
Door toggles position &
PASS \\
\hline
\end{tabular}
\end{table}
\FloatBarrier

All eight predefined functional scenarios
were reported as passing during prototype
testing.

This result indicates that the observed
behavior matched the expected outcomes
for the selected scenarios.

However, functional test completion
should not be interpreted as a measured
100\% fire-detection accuracy.

A statistically meaningful accuracy
evaluation would require additional
labeled fire and non-fire samples,
repeated trials, and analysis of
false-positive and false-negative events.

\subsection{Prototype Demonstration}

The physical prototype integrates
sensors, warning indicators, an LCD,
and the SG90 door mechanism.

A photograph of the completed hardware
will be inserted in the final report
after the demonstration images
have been collected.

% ----- FIGURE 4: PROTOTYPE PHOTO -----
\begin{figure}[!htbp]
    \centering
    \IfFileExists{src/prototype.png}{
        \includegraphics[
            width=0.80\linewidth,
            height=5cm,
            keepaspectratio
        ]{src/prototype.png}
    }{
        \fbox{
            \parbox[c][2.5cm][c]{0.80\linewidth}{
                \centering
                PHYSICAL PROTOTYPE PHOTO\\
                TO BE INSERTED
            }
        }
    }
    \caption{Physical prototype of the
    ESP32-based fire detection system.}
    \label{fig:prototype}
\end{figure}
\FloatBarrier

\subsection{Discussion}

The experimental observations support
the basic functionality of the
three-state warning mechanism and
servo-control logic.

One advantage of the proposed design
is that warning decisions are processed
directly on the ESP32. Consequently,
the local alarm does not require
Internet connectivity.

The prioritized DANGER condition
also allows the system to respond
when a single sensor reports
a potentially hazardous event.

However, the same design may trigger
unwanted alarms when smoke or
gas detection is caused by
non-fire activities.

Research on multisensor fusion and
signal analysis has demonstrated
alternative approaches to distinguishing
fire conditions from environmental
interference
\cite{yu2023multidetector,
solorzano2022gas,nakip2021rtpnn}.

The present prototype does not
implement those advanced classification
techniques. Instead, it prioritizes
simplicity and functional demonstration.

Additional limitations include
the absence of independently measured
response latency, sensor accuracy,
false-alarm rate, and long-term
reliability data.

The assembled LM35 supply configuration
and the interface voltage levels of
5 V modules also require hardware-level
verification.

Furthermore, the servo represents
a small-scale demonstration mechanism
rather than a certified emergency
door-control system.

% ==========================================
\section{Conclusion}
% ==========================================

This project developed an ESP32-based
fire detection and warning prototype
using LM35, MQ-2, and flame sensors.

A rule-based algorithm classifies
environmental conditions into
SAFE, WARNING, and DANGER states.

The system provides local visual
and audible warnings through three
LEDs, an LCD1602 display, and
an active buzzer.

An SG90 servo motor demonstrates
automatic door opening during
dangerous conditions, while
a push button supports manual
control outside the DANGER state.

The group reported successful
functional behavior across eight
predefined test scenarios.

The prototype demonstrates the
feasibility of integrating multiple
sensors, local warning indicators,
and actuator control using ESP32.
It is an educational prototype
and has not been validated as a
certified fire-safety device.

Future work may include sensor
calibration, response-time measurement,
false-alarm evaluation, adaptive
warning policies, and remote
monitoring through platforms such
as Blynk.

These enhancements should preserve
local safety functions even if
network connectivity is unavailable.

% ==========================================
\section*{Author Contributions}
% ==========================================

All three members contributed equally
to the project, with complementary
responsibilities.

\textbf{Hoang Nhat Tam (SE202018):}
Project coordination, task allocation,
literature research, report preparation,
and design review.

\textbf{Phan Tan Hung (SE201027):}
ESP32 firmware development,
sensor integration, warning logic,
debugging, and functional testing.

\textbf{Le Vo Tuan Kiet (SE203376):}
Hardware assembly, wiring,
prototype integration, actuator testing,
and hardware troubleshooting.

All members participated in system
review, functional evaluation,
and demonstration preparation.

% ==========================================
\begin{thebibliography}{99}
% ==========================================

\bibitem{yu2023multidetector}
M. Yu, H. Yuan, K. Li, and J. Wang,
``Research on multi-detector real-time fire
alarm technology based on signal similarity,''
\textit{Fire Safety Journal},
vol. 136, Art. no. 103724, 2023.
doi: 10.1016/j.firesaf.2022.103724.

\bibitem{fathoni2025multisensor}
K. Fathoni, A. F. Suni, N. Iksan,
and M. S. Aziz,
``Multisensor Integration in Early Fire and Gas
Detection Based on Artificial Neural Network,''
\textit{Edu Komputika Journal},
vol. 12, no. 2, pp. 157--167, 2025.
doi: 10.15294/edukom.v12i2.27436.

\bibitem{purba2026smartsafety}
J. Purba, D. Kiswanto, J. B. Henrydunan,
and R. Dly,
``Smart Safety Room: ESP32 Decision Tree-Based
Multi-Hazard Detection System,''
\textit{Journal of Artificial Intelligence
and Engineering Applications},
vol. 5, no. 2, pp. 2627--2637, 2026.
doi: 10.59934/jaiea.v5i2.1947.

\bibitem{ramadani2026blynk}
N. Ramadani,
``IoT-Based Early Fire Detection System using
Flame and Gas Sensors with Telegram Bot and
Blynk Dashboard Integration,''
\textit{Journal of Artificial Intelligence
and Engineering Applications},
vol. 6, no. 1, pp. 603--612, 2026.
doi: 10.59934/jaiea.v6i1.2805.

\bibitem{solorzano2022gas}
A. Solorzano et al.,
``Early fire detection based on gas sensor
arrays: Multivariate calibration and validation,''
\textit{Sensors and Actuators B: Chemical}, 2022.
doi: 10.1016/j.snb.2021.130961.

\bibitem{nakip2021rtpnn}
M. Nakip, C. Guzelis, and O. Yildiz,
``Recurrent Trend Predictive Neural Network
for Multi-Sensor Fire Detection,''
\textit{IEEE Access},
vol. 9, pp. 84204--84216, 2021.
doi: 10.1109/ACCESS.2021.3087736.

\bibitem{rehman2021smartfire}
A. Rehman et al.,
``Smart Fire Detection and Deterrent System
for Human Savior by Using Internet of Things
(IoT),''
\textit{Energies},
vol. 14, no. 17, Art. no. 5500, 2021.
doi: 10.3390/en14175500.

\bibitem{espressif_adc_calibration}
Espressif Systems,
``Analog to Digital Converter (ADC)
Calibration Driver,''
\textit{ESP-IDF Programming Guide}.
[Online]. Available:
\texttt{https://docs.espressif.com/projects/esp-idf/}

\end{thebibliography}

% ==========================================
\section*{AI Usage Declaration}
% ==========================================

Generative Artificial Intelligence (AI),
particularly ChatGPT by OpenAI, was used
as a supporting tool during the development
and documentation of this project.

AI assistance was primarily used for
exploring scientific literature,
understanding technical concepts,
reviewing system architecture,
discussing sensor integration and
fire detection logic, and improving
the organization and academic language
of the final report.

The project team critically reviewed
AI-generated suggestions against
the implemented source code,
hardware design, technical documentation,
and available functional testing results.
AI-generated content was not automatically
accepted as technically correct.

In particular, the team distinguished
between implemented features and
proposed future improvements, including
remote monitoring through IoT platforms
such as Blynk.

All final design decisions,
implementation activities,
experimental observations,
and conclusions remain the
responsibility of the project members.

Selected AI-assisted decisions,
verification activities, and reflections
are documented separately in the
AI Audit Log.

\end{document}
